\chapter{Analyse et conception}

Ce chapitre présente la modélisation du système au moyen de diagrammes UML : vue fonctionnelle (cas d'utilisation), vue structurelle (classes), vue dynamique (activités et états-transitions). Les diagrammes eux-mêmes sont fournis dans les annexes graphiques ; ce chapitre en donne les descriptions textuelles et les commentaires.

\section{Diagramme de cas d'utilisation}

Le diagramme de cas d'utilisation du système met en relation trois acteurs (rédacteur, visiteur téléchargeur, administrateur) et les cas d'utilisation suivants :

\begin{itemize}
    \item Pour le rédacteur : s'identifier, créer un document, inviter des collaborateurs, modifier un document, consulter les versions, restaurer une version, créer un lien de téléchargement.
    \item Pour le visiteur : télécharger un journal, consulter l'état de la demande de téléchargement.
    \item Pour l'administrateur : consulter les documents ouverts, consulter les liens de téléchargement, consulter les rédacteurs bloqués.
\end{itemize}

\begin{infobox}
Le cas \textit{Créer un lien de téléchargement} inclut la consultation des versions pour garantir qu'un état cohérent est référencé. Le cas \textit{Télécharger un journal} inclut la consultation de l'état de la demande pour informer le visiteur. Le cas \textit{Restaurer une version} étend \textit{Consulter les versions}.
\end{infobox}

\subsection{Descriptions textuelles des cas d'utilisation}

Les descriptions ci-dessous reprennent, pour chaque cas d'utilisation identifié, l'acteur principal, les préconditions, le scénario nominal, les scénarios alternatifs ou exceptions, et les postconditions.

\begin{description}[leftmargin=0pt, style=nextline, font=\normalfont\bfseries\color{violetprincipal}]

    \item[UC01 --- S'identifier]

    \textbf{Acteur principal :} Rédacteur.

    \textbf{Objectif :} Permettre à un rédacteur de fournir son identité afin d'utiliser les fonctionnalités de l'application.

    \textbf{Préconditions :} L'application est accessible.

    \textbf{Scénario nominal :}
    \begin{enumerate}
        \item Le rédacteur accède à l'application.
        \item Le système demande son nom et son adresse e-mail.
        \item Le rédacteur renseigne son nom et son adresse e-mail.
        \item Le système vérifie la validité des informations fournies.
        \item Le système associe l'identité fournie au rédacteur concerné, selon le mécanisme d'identification retenu.
        \item Le système permet au rédacteur d'accéder aux fonctionnalités autorisées.
    \end{enumerate}

    \textbf{Scénarios alternatifs et exceptions :}
    \begin{itemize}
        \item \textbf{A1} : Le nom ou l'adresse e-mail est manquant. Le système demande de compléter les informations.
        \item \textbf{A2} : L'adresse e-mail est invalide. Le système demande une correction.
        \item \textbf{A3} : Une erreur survient pendant l'identification. Le système informe le rédacteur.
    \end{itemize}

    \textbf{Postconditions :} L'identité déclarée du rédacteur est disponible ; le rédacteur peut accéder aux fonctionnalités correspondant à ses droits.

    \textit{Remarque :} Le sujet précise qu'il n'y a pas de véritables comptes rédacteurs. Ce cas décrit donc la fourniture et la reconnaissance de l'identité déclarée, et non une authentification par mot de passe.

    \item[UC02 --- Créer un document]

    \textbf{Acteur principal :} Rédacteur.

    \textbf{Objectif :} Permettre à un rédacteur de créer un nouveau document textuel dans l'application.

    \textbf{Préconditions :}
    \begin{itemize}
        \item Le rédacteur a renseigné son nom et son adresse e-mail.
        \item Le rédacteur n'est pas temporairement bloqué.
    \end{itemize}

    \textbf{Scénario nominal :}
    \begin{enumerate}
        \item Le rédacteur demande la création d'un document.
        \item Le système vérifie que le rédacteur peut effectuer cette opération.
        \item Le système vérifie que le rédacteur n'est pas bloqué.
        \item Le système crée un nouveau document vide.
        \item Le système associe le document à son créateur.
        \item Le système initialise les informations du document.
        \item Le système affiche le document dans l'éditeur.
        \item Le rédacteur peut commencer à écrire.
    \end{enumerate}

    \textbf{Scénarios alternatifs et exceptions :}
    \begin{itemize}
        \item \textbf{A1} : Le rédacteur est temporairement bloqué. Le système refuse la création si cette opération est interdite aux rédacteurs bloqués.
        \item \textbf{A2} : Une erreur survient pendant la création. Le système informe le rédacteur de l'échec.
    \end{itemize}

    \textbf{Postconditions :} Le document est créé et associé à son créateur ; il est disponible pour les opérations autorisées.

    \item[UC03 --- Inviter des collaborateurs]

    \textbf{Acteur principal :} Rédacteur.

    \textbf{Objectif :} Permettre au propriétaire d'un document de donner accès à d'autres rédacteurs afin qu'ils puissent collaborer.

    \textbf{Préconditions :}
    \begin{itemize}
        \item Le document existe.
        \item Le rédacteur est le propriétaire du document.
        \item Le rédacteur n'est pas bloqué, si le blocage interdit cette opération.
    \end{itemize}

    \textbf{Scénario nominal :}
    \begin{enumerate}
        \item Le rédacteur ouvre le document concerné.
        \item Le rédacteur sélectionne la fonctionnalité de partage.
        \item Le système vérifie les droits du rédacteur.
        \item Le rédacteur renseigne les informations nécessaires pour identifier le collaborateur, notamment son adresse e-mail.
        \item Le système vérifie les informations fournies.
        \item Le système enregistre l'accès du collaborateur au document.
        \item Le système confirme le partage au rédacteur.
    \end{enumerate}

    \textbf{Scénarios alternatifs et exceptions :}
    \begin{itemize}
        \item \textbf{A1} : Le rédacteur n'est pas propriétaire du document. Le système refuse l'opération.
        \item \textbf{A2} : Les informations du collaborateur sont invalides. Le système demande leur correction.
        \item \textbf{A3} : Le collaborateur possède déjà un accès. Le système informe le rédacteur.
        \item \textbf{A4} : Une erreur survient pendant le partage. Le système signale l'échec.
    \end{itemize}

    \textbf{Postconditions :} Le collaborateur dispose d'un accès au document conformément aux droits définis ; le système conserve les informations nécessaires à la gestion du partage.

    \item[UC04 --- Modifier un document]

    \textbf{Acteur principal :} Rédacteur.

    \textbf{Acteurs secondaires :} Autres rédacteurs collaborateurs.

    \textbf{Objectif :} Permettre à un rédacteur autorisé de modifier un document, éventuellement en collaboration avec d'autres rédacteurs.

    \textbf{Préconditions :}
    \begin{itemize}
        \item Le rédacteur est identifié.
        \item Le document existe.
        \item Le rédacteur est propriétaire du document ou dispose d'un accès collaborateur.
        \item Le rédacteur n'est pas bloqué.
    \end{itemize}

    \textbf{Scénario nominal :}
    \begin{enumerate}
        \item Le rédacteur ouvre le document.
        \item Le système vérifie ses droits d'accès.
        \item Le système vérifie que le rédacteur peut écrire.
        \item Le système vérifie que son crédit d'écriture est disponible.
        \item Le système affiche le contenu actuel du document.
        \item Le système indique les collaborateurs présents et leurs positions d'écriture.
        \item Le rédacteur saisit ou modifie du texte.
        \item Le système contrôle la quantité de caractères que le rédacteur peut encore écrire.
        \item Le système accepte les modifications respectant les règles de crédit.
        \item Le système applique les modifications sans perdre celles des autres collaborateurs.
        \item Le système propage les modifications aux collaborateurs en quasi-temps réel.
        \item Le système met à jour le crédit d'écriture du rédacteur.
        \item Le système conserve les modifications acceptées et l'historique nécessaire.
    \end{enumerate}

    \textbf{Scénarios alternatifs et exceptions :}
    \begin{itemize}
        \item \textbf{A1} : Le rédacteur ne possède pas les droits nécessaires. Le système refuse l'accès en écriture.
        \item \textbf{A2} : Le rédacteur est bloqué. Le système lui interdit d'écrire.
        \item \textbf{A3} : Le crédit est épuisé. Le système refuse les nouvelles écritures jusqu'à la recharge du crédit.
        \item \textbf{A4} : La saisie dépasse le crédit restant. Le système limite ou refuse la partie excédentaire, selon la stratégie d'édition retenue.
        \item \textbf{A5} : Plusieurs rédacteurs modifient simultanément le document. Le système coordonne les modifications afin de préserver toutes les modifications acceptées.
        \item \textbf{A6} : Une tentative de contournement des limites est détectée. Le système applique le mécanisme de détection et de blocage temporaire prévu.
        \item \textbf{A7} : Une erreur de connexion survient. Le système informe le rédacteur et applique la stratégie de récupération prévue.
    \end{itemize}

    \textbf{Postconditions :}
    \begin{itemize}
        \item Les modifications acceptées sont conservées.
        \item Les crédits sont mis à jour sans dépassement.
        \item Les collaborateurs convergent vers un même état du document.
    \end{itemize}

    \item[UC05 --- Consulter les versions]

    \textbf{Acteur principal :} Rédacteur.

    \textbf{Objectif :} Permettre à un rédacteur autorisé de consulter l'historique des versions d'un document.

    \textbf{Préconditions :}
    \begin{itemize}
        \item Le document existe.
        \item Le rédacteur dispose des droits nécessaires pour consulter le document.
    \end{itemize}

    \textbf{Scénario nominal :}
    \begin{enumerate}
        \item Le rédacteur sélectionne le document concerné.
        \item Le rédacteur demande à consulter les versions.
        \item Le système vérifie les droits d'accès du rédacteur.
        \item Le système récupère l'historique des versions disponibles.
        \item Le système affiche la liste des versions avec leurs informations pertinentes : numéro, date et auteur, selon les données conservées.
        \item Le rédacteur sélectionne une version.
        \item Le système affiche le contenu de la version sélectionnée.
        \item Le rédacteur consulte la version.
    \end{enumerate}

    \textbf{Scénarios alternatifs et exceptions :}
    \begin{itemize}
        \item \textbf{A1} : Le document ne possède pas de version historique disponible. Le système informe le rédacteur.
        \item \textbf{A2} : Le rédacteur ne dispose pas des droits nécessaires. Le système refuse l'accès.
        \item \textbf{A3} : La version sélectionnée est indisponible. Le système informe le rédacteur.
        \item \textbf{A4} : Une erreur survient lors du chargement de l'historique. Le système signale l'échec.
    \end{itemize}

    \textbf{Postconditions :} Le rédacteur a pu consulter les versions accessibles ; le contenu courant n'est pas modifié par la simple consultation.

    \textit{Remarque :} La restauration d'une version constitue un cas d'utilisation distinct (UC06).

    \item[UC06 --- Restaurer une version]

    \textbf{Acteur principal :} Rédacteur.

    \textbf{Objectif :} Permettre à un rédacteur autorisé de restaurer une version antérieure d'un document.

    \textbf{Préconditions :}
    \begin{itemize}
        \item Le document existe.
        \item Le rédacteur dispose des droits nécessaires.
        \item Le document possède au moins une version antérieure consultable.
    \end{itemize}

    \textbf{Scénario nominal :}
    \begin{enumerate}
        \item Le rédacteur ouvre le document.
        \item Le rédacteur demande la consultation de l'historique.
        \item Le système affiche les versions disponibles.
        \item Le rédacteur sélectionne une version.
        \item Le système affiche le contenu de la version sélectionnée.
        \item Le rédacteur demande la restauration.
        \item Le système vérifie les droits du rédacteur.
        \item Le système vérifie les conditions de restauration.
        \item Le système restaure le contenu correspondant à la version sélectionnée.
        \item Le système conserve l'historique et enregistre cette restauration.
        \item Le système propage le nouvel état du document aux collaborateurs.
    \end{enumerate}

    \textbf{Scénarios alternatifs et exceptions :}
    \begin{itemize}
        \item \textbf{A1} : Aucune version antérieure n'est disponible. Le système informe le rédacteur.
        \item \textbf{A2} : La version sélectionnée n'existe plus ou n'est pas accessible. Le système refuse la restauration.
        \item \textbf{A3} : Le rédacteur ne dispose pas des droits nécessaires. Le système refuse l'opération.
        \item \textbf{A4} : Des modifications simultanées sont en cours. Le système coordonne la restauration avec les autres modifications pour éviter toute perte de données.
        \item \textbf{A5} : Une erreur survient pendant la restauration. Le système signale l'échec sans présenter une restauration incomplète comme réussie.
    \end{itemize}

    \textbf{Postconditions :} Le contenu du document correspond à la version restaurée, sous réserve de la stratégie de restauration retenue ; l'historique est préservé ; les collaborateurs peuvent consulter l'état restauré.

    \item[UC07 --- Créer un lien de téléchargement]

    \textbf{Acteur principal :} Rédacteur.

    \textbf{Objectif :} Permettre à un rédacteur autorisé de créer un lien de téléchargement pour un document, avec un quota défini.

    \textbf{Préconditions :}
    \begin{itemize}
        \item Le document existe.
        \item Le rédacteur dispose des droits nécessaires.
        \item Le rédacteur n'est pas bloqué, si le blocage interdit cette opération.
    \end{itemize}

    \textbf{Scénario nominal :}
    \begin{enumerate}
        \item Le rédacteur ouvre le document concerné.
        \item Le rédacteur demande la création d'un lien de téléchargement.
        \item Le système vérifie ses droits.
        \item Le système demande le nombre de téléchargements autorisés.
        \item Le rédacteur saisit le quota souhaité.
        \item Le système vérifie que le quota est valide et strictement positif.
        \item Le système crée un lien unique associé au document.
        \item Le système initialise le quota total et le nombre de téléchargements restants.
        \item Le système active le lien.
        \item Le système communique le lien au rédacteur.
    \end{enumerate}

    \textbf{Scénarios alternatifs et exceptions :}
    \begin{itemize}
        \item \textbf{A1} : Le quota fourni est invalide. Le système demande une nouvelle valeur.
        \item \textbf{A2} : Le rédacteur ne dispose pas des droits nécessaires. Le système refuse la création.
        \item \textbf{A3} : Une erreur survient pendant la création. Le système informe le rédacteur de l'échec.
    \end{itemize}

    \textbf{Postconditions :} Un lien de téléchargement est créé, associé au document, avec quota initialisé ; il permet de télécharger le document dans la limite autorisée.

    \item[UC08 --- Télécharger un journal]

    \textbf{Acteur principal :} Visiteur.

    \textbf{Objectif :} Permettre à un visiteur anonyme de télécharger le document au format PDF, dans la limite du quota du lien.

    \textbf{Préconditions :}
    \begin{itemize}
        \item Le visiteur possède un lien de téléchargement.
        \item Le lien correspond à un lien reconnu par le système.
    \end{itemize}

    \textbf{Scénario nominal :}
    \begin{enumerate}
        \item Le visiteur ouvre le lien de téléchargement.
        \item Le système vérifie l'existence et la validité du lien.
        \item Le système vérifie l'état du lien et son quota disponible.
        \item Le système vérifie qu'une place de téléchargement peut être attribuée.
        \item Le système réserve atomiquement une place dans le quota afin de prévenir tout dépassement en cas de demandes simultanées.
        \item Le système indique au visiteur que sa demande est en cours de traitement.
        \item Le système obtient un état cohérent du document au moment de la prise en charge du téléchargement.
        \item Le système génère le PDF correspondant à cet état.
        \item Le système met le fichier à disposition du visiteur.
        \item Le système indique que le téléchargement est terminé ou fournit les informations nécessaires à sa récupération.
        \item Le système finalise le traitement du téléchargement.
        \item Si le dernier téléchargement autorisé est consommé, le système informe le créateur du lien par e-mail.
    \end{enumerate}

    \textbf{Scénarios alternatifs et exceptions :}
    \begin{itemize}
        \item \textbf{A1} : Le lien est invalide. Le système informe le visiteur.
        \item \textbf{A2} : Le quota est épuisé. Le système indique clairement que le journal n'est plus disponible.
        \item \textbf{A3} : Le visiteur arrive pendant un pic de charge. Le système lui indique si son téléchargement est en cours, en attente ou si le quota est épuisé.
        \item \textbf{A4} : Le quota est consommé pendant le traitement d'autres demandes. Le système refuse les nouvelles demandes non réservées.
        \item \textbf{A5} : Le PDF ne peut pas être généré. Le système applique une politique de traitement des échecs cohérente avec la réservation du quota.
        \item \textbf{A6} : Le téléchargement échoue pendant la transmission. Le système applique la politique prévue pour le quota et informe le visiteur.
        \item \textbf{A7} : Le dernier exemplaire vient d'être distribué. Le système informe rapidement les visiteurs encore en attente que le quota est épuisé.
    \end{itemize}

    \textbf{Postconditions :}
    \begin{itemize}
        \item Le téléchargement autorisé est traité conformément aux règles de quota.
        \item Le quota ne dépasse jamais le nombre autorisé.
        \item Le visiteur reçoit une information sur l'état de sa demande.
        \item Le document téléchargé correspond à un état complet et cohérent.
    \end{itemize}

    \item[UC09 --- Consulter l'état de la demande de téléchargement]

    \textbf{Acteur principal :} Visiteur.

    \textbf{Objectif :} Permettre au visiteur de connaître l'état de sa demande de téléchargement.

    \textbf{Préconditions :}
    \begin{itemize}
        \item Le visiteur a initié une demande de téléchargement.
        \item Le système dispose d'un moyen permettant de retrouver cette demande, par exemple un identifiant ou un mécanisme de suivi.
    \end{itemize}

    \textbf{Scénario nominal :}
    \begin{enumerate}
        \item Le visiteur accède à l'interface de suivi de sa demande.
        \item Le visiteur fournit ou utilise les informations de suivi disponibles.
        \item Le système recherche la demande correspondante.
        \item Le système détermine l'état actuel de la demande.
        \item Le système affiche l'état au visiteur.
    \end{enumerate}

    \textbf{Scénarios alternatifs et exceptions :}
    \begin{itemize}
        \item \textbf{A1} : La demande est en attente. Le système indique que le visiteur doit patienter.
        \item \textbf{A2} : La demande est en cours de traitement. Le système informe le visiteur.
        \item \textbf{A3} : La demande est terminée. Le système indique que le téléchargement est disponible ou terminé.
        \item \textbf{A4} : Le quota est épuisé avant la prise en charge. Le système informe le visiteur que le téléchargement n'est plus disponible.
        \item \textbf{A5} : La demande est refusée. Le système indique le motif lorsque celui-ci peut être communiqué.
        \item \textbf{A6} : La demande est introuvable ou l'identifiant de suivi est invalide. Le système informe le visiteur.
    \end{itemize}

    \textbf{Postconditions :} Le visiteur connaît l'état de sa demande ; aucune modification du quota n'est effectuée par la simple consultation.

    \item[UC10 --- Consulter les documents ouverts]

    \textbf{Acteur principal :} Administrateur.

    \textbf{Objectif :} Permettre à l'administrateur de superviser les documents actuellement ouverts dans l'application.

    \textbf{Préconditions :}
    \begin{itemize}
        \item L'administrateur est identifié.
        \item L'administrateur dispose des droits de supervision.
    \end{itemize}

    \textbf{Scénario nominal :}
    \begin{enumerate}
        \item L'administrateur accède à l'espace de supervision.
        \item Il sélectionne la fonctionnalité de consultation des documents ouverts.
        \item Le système vérifie ses droits.
        \item Le système recherche les documents actuellement ouverts.
        \item Le système récupère les informations de supervision disponibles.
        \item Le système affiche les documents ouverts.
        \item L'administrateur consulte les informations affichées, notamment les documents concernés et, si disponible, les rédacteurs présents.
        \item L'administrateur termine la consultation.
    \end{enumerate}

    \textbf{Scénarios alternatifs et exceptions :}
    \begin{itemize}
        \item \textbf{A1} : Aucun document n'est actuellement ouvert. Le système affiche un message approprié.
        \item \textbf{A2} : L'administrateur ne possède pas les droits nécessaires. Le système refuse l'accès.
        \item \textbf{A3} : Un document se ferme pendant la consultation. Le système actualise son état selon le mécanisme de supervision retenu.
        \item \textbf{A4} : Une erreur survient lors de la récupération des données. Le système informe l'administrateur.
    \end{itemize}

    \textbf{Postconditions :} L'administrateur a pu consulter l'état des documents ouverts ; aucune modification du contenu n'est effectuée par la simple consultation.

    \textit{Remarque :} Les informations de présence doivent correspondre à l'état réel des sessions. Le système doit prévoir comment traiter les sessions interrompues ou les déconnexions non signalées.

    \item[UC11 --- Consulter les liens de téléchargement]

    \textbf{Acteur principal :} Administrateur.

    \textbf{Objectif :} Permettre à l'administrateur de consulter les liens de téléchargement actifs et leurs quotas restants.

    \textbf{Préconditions :}
    \begin{itemize}
        \item L'administrateur est identifié.
        \item L'administrateur dispose des droits de supervision.
    \end{itemize}

    \textbf{Scénario nominal :}
    \begin{enumerate}
        \item L'administrateur accède à l'espace de supervision.
        \item Il sélectionne la fonctionnalité de consultation des liens de téléchargement.
        \item Le système vérifie ses droits.
        \item Le système récupère les liens de téléchargement à superviser.
        \item Le système récupère les informations associées à chaque lien.
        \item Le système affiche les informations pertinentes, notamment : le document concerné, l'état du lien, le quota initial, le nombre de téléchargements consommés, le nombre de téléchargements restants.
        \item L'administrateur consulte les informations.
    \end{enumerate}

    \textbf{Scénarios alternatifs et exceptions :}
    \begin{itemize}
        \item \textbf{A1} : Aucun lien n'est disponible. Le système affiche un message approprié.
        \item \textbf{A2} : L'administrateur ne possède pas les droits nécessaires. Le système refuse l'accès.
        \item \textbf{A3} : Certains liens sont épuisés ou désactivés. Le système affiche leur état respectif.
        \item \textbf{A4} : Une erreur survient lors du chargement des données. Le système informe l'administrateur.
    \end{itemize}

    \textbf{Postconditions :} L'administrateur a pu consulter les informations disponibles sur les liens ; aucune modification du quota ou de l'état des liens n'est effectuée par la simple consultation.

    \textit{Remarque :} La consultation porte sur les liens et leurs quotas. Une opération de modification ou de désactivation des liens nécessiterait un cas d'utilisation supplémentaire si elle est retenue.

    \item[UC12 --- Consulter les rédacteurs bloqués]

    \textbf{Acteur principal :} Administrateur.

    \textbf{Objectif :} Permettre à l'administrateur de consulter la liste des rédacteurs temporairement bloqués par le système.

    \textbf{Préconditions :}
    \begin{itemize}
        \item L'administrateur est identifié et autorisé à accéder à la fonctionnalité de supervision.
    \end{itemize}

    \textbf{Scénario nominal :}
    \begin{enumerate}
        \item L'administrateur accède à l'espace de supervision.
        \item Il sélectionne la fonctionnalité de consultation des rédacteurs bloqués.
        \item Le système vérifie ses droits.
        \item Le système recherche les rédacteurs actuellement bloqués.
        \item Le système récupère les informations nécessaires à leur identification et à leur suivi.
        \item Le système affiche la liste des rédacteurs bloqués.
        \item L'administrateur consulte les informations disponibles.
    \end{enumerate}

    \textbf{Scénarios alternatifs et exceptions :}
    \begin{itemize}
        \item \textbf{A1} : Aucun rédacteur n'est actuellement bloqué. Le système affiche une liste vide ou un message approprié.
        \item \textbf{A2} : L'administrateur ne possède pas les droits nécessaires. Le système refuse l'accès.
        \item \textbf{A3} : Une erreur survient lors de la récupération des informations. Le système informe l'administrateur.
    \end{itemize}

    \textbf{Postconditions :} L'administrateur a pu consulter la liste des rédacteurs bloqués ; aucune modification du statut des rédacteurs n'est effectuée par la simple consultation.

    \textit{Remarque :} Le déblocage d'un rédacteur n'est pas décrit comme un cas d'utilisation distinct dans le diagramme présenté. Il ne faut donc pas le considérer comme une fonctionnalité acquise sans validation.
\end{description}

\section{Diagramme de classes}

Le diagramme de classes du système comprend les entités principales suivantes et leurs relations :

\begin{itemize}
    \item \textbf{Rédacteur} : id, nom, email, crédit restant ; participe à un ou plusieurs documents ; possède un quota d'écriture par jour ; peut être bloqué.
    \item \textbf{Document} : id, titre, contenu, version active ; possède un ensemble de versions (au moins une) ; peut générer plusieurs liens de téléchargement.
    \item \textbf{Version} : id, numéro, contenu, horodatage ; appartient à un document ; peut être référencée par un téléchargement.
    \item \textbf{LienTéléchargement} : id, quota, téléchargements consommés, créateur ; associé à un document ; concerne plusieurs téléchargements.
    \item \textbf{Téléchargement} : id, lien, version référencée, état, date de demande ; traite la génération du PDF.
    \item \textbf{QuotaÉcriture} : rédacteur, date, caractères utilisés, limite ; consommé par les modifications.
    \item \textbf{Blocage} : id, rédacteur, motif, dates de début et de fin.
\end{itemize}

\begin{infobox}
La multiplicité \texttt{1..*} entre \texttt{Document} et \texttt{Version} traduit le fait qu'un document possède au moins une version et en accumule à chaque sauvegarde. Le lien de téléchargement n'est pas lié à une version figée : c'est la \texttt{Téléchargement} qui référence la version au moment de la demande, garantissant le principe de validité du lien.
\end{infobox}

\section{Diagrammes d'activités}

\subsection{Distribution d'un document via un lien}

Le processus de distribution d'un document via un lien se déroule ainsi :

\begin{enumerate}
    \item Le visiteur demande le téléchargement.
    \item Le système vérifie si le quota est disponible. Si non, il refuse avec la mention « quota épuisé ».
    \item Si oui, le système réserve atomiquement une place.
    \item Le système fige l'état du document (référence à une version).
    \item Le système génère le PDF.
    \item Si la génération réussit, le système livre le PDF au visiteur. Si la génération échoue, la place est rendue et le visiteur est notifié.
    \item Après livraison, le système vérifie si la place consommée était la dernière. Si oui, il notifie le créateur par e-mail.
\end{enumerate}

\begin{infobox}
La place est réservée \emph{avant} la génération du PDF (réservation atomique). Si la génération échoue, la place est rendue, ce qui garantit qu'un lien à quota $N$ permet bien $N$ téléchargements réussis, même en cas de panne transitoire, sans jamais dépasser $N$.
\end{infobox}

\subsection{Modification d'un document}

Le processus de modification d'un document par un rédacteur se déroule ainsi :

\begin{enumerate}
    \item Le rédacteur ouvre le document.
    \item Le système vérifie ses droits d'écriture. Si non, il refuse.
    \item Le système vérifie que le rédacteur n'est pas bloqué. Si oui, il refuse.
    \item Le système vérifie que le crédit est disponible. Si non, il refuse les nouvelles écritures.
    \item Le rédacteur saisit ou modifie du texte.
    \item Le système fusionne sans perte avec l'état courant.
    \item Le système propage la modification aux collaborateurs.
    \item Le système met à jour le crédit du rédacteur.
    \item Le système conserve la modification et l'historique.
\end{enumerate}

\section{Diagrammes d'états-transitions}

\subsection{Cycle de vie d'une demande de téléchargement}
\label{sec:cycle-vie-demande}

Une demande de téléchargement passe par les états suivants :

\begin{itemize}
    \item \textbf{Reçue} : la demande est enregistrée, avant décision.
    \item \textbf{En attente} : la demande est placée en file si la capacité ou les places disponibles ne permettent pas une admission immédiate.
    \item \textbf{Autorisée} : une place a été réservée atomiquement. C'est à cette transition que la place est consommée au sens de la réservation.
    \item \textbf{Refusée « épuisé »} : aucune place n'est disponible. Cet état est définitif pour un lien donné.
    \item \textbf{PDF prêt} : le PDF a été généré pour l'état figé du document.
    \item \textbf{Terminée} : le PDF a été livré. La place est définitivement consommée au sens de la livraison.
    \item \textbf{Échouée} : la génération du PDF a échoué ou le visiteur n'a pas abouti. La place est rendue.
\end{itemize}

\begin{infobox}
La place est \textbf{réservée} à la transition \textit{Reçue} $\to$ \textit{Autorisée}, ce qui empêche tout dépassement en cas de concurrence. Elle est \textbf{livrée} à la transition \textit{PDF prêt} $\to$ \textit{Terminée}, et \textbf{rendue} à la transition \textit{PDF prêt} $\to$ \textit{Échouée}. Aucune demande ne peut passer de \textit{Refusée} à un autre état : le quota épuisé est définitif pour un lien donné.
\end{infobox}

\subsection{Cycle de vie d'un document et des versions}

Un document passe par les états suivants : \textit{Créé}, \textit{En édition}, \textit{Version sauvegardée}, \textit{Version restaurée}. Les transitions sont :

\begin{itemize}
    \item \textit{Créé} $\to$ \textit{En édition} : à l'ouverture pour édition.
    \item \textit{En édition} $\to$ \textit{Version sauvegardée} : à chaque sauvegarde explicite ou fin de session.
    \item \textit{Version sauvegardée} $\to$ \textit{En édition} : à la reprise d'édition.
    \item \textit{En édition} $\to$ \textit{Version restaurée} : à la demande de restauration.
    \item \textit{Version restaurée} $\to$ \textit{Version sauvegardée} : la restauration crée une nouvelle version.
\end{itemize}

\begin{infobox}
Restaurer ne détruit pas l'historique : cela crée une nouvelle version dont le contenu est celui de la version restaurée (cf. section~\ref{sec:auth-admin} et les hypothèses sur la gestion des versions).
\end{infobox}

% ============================================
